\input{preamble.tex}

\title{Отчет по лабораторной работе № 11. \\Настройка NAT. Планирование}
\author{Сущенко Алина Николаевна \\ Группа: НПИбд-01-23}
\date{2026}

\begin{document}

\maketitle
\newpage

\tableofcontents

\newpage
\section{Цель работы}
Провести подготовительные мероприятия по подключению локальной сети организации к Интернету: построить схему соединения локальной сети с сетью провайдера и модельной сетью Интернета, настроить IP-адреса и DNS-записи.

\newpage
\section{Выполнение работы}

\subsection{Построение схемы L1 (физический уровень)}

Перед началом настройки сети необходимо разместить оборудование провайдера и модельного Интернета в физической рабочей области Packet Tracer. (Рис. \ref{01}):

\begin{center}
    \centering
    \includegraphics[width=\textwidth]{../images/01.png}
    \captionof{figure}{Обновленная схема топологии сети.}
    \label{01}
\end{center}

\begin{enumerate}
\item В физической рабочей области добавлены здание провайдера и здание, имитирующее расположение серверов модельного Интернета (Рис. \ref{02}):

\begin{center}
    \centering
    \includegraphics[width=\textwidth]{../images/02.png}
    \captionof{figure}{Схема сети в физической рабочей области Packet Tracer.}
    \label{02}
\end{center}

\item Перенесем из Донской оборудование в здание сети провайдера (Рис. \ref{03}):

\begin{center}
    \centering
    \includegraphics[width=\textwidth]{../images/03.png}
    \captionof{figure}{Оборудование в здании сети провайдера.}
    \label{03}
\end{center}

\item Так же перенесем в физическом окружении оборудование в здании сети модельного Интернета: коммутатор \texttt{internet-sw-1}, серверы \texttt{www.rudn.ru}, \texttt{stud.rudn.university}, \texttt{www.yandex.ru}, \texttt{system.pf.ur.ru}, а также медиаконвертер \texttt{ansusenko-mc-01}.} (Рис. \ref{04}):

\begin{center}
    \centering
    \includegraphics[width=\textwidth]{../images/04.png}
    \captionof{figure}{Оборудование в здании сети модельного Интернета.}
    \label{04}
\end{center}

\item Настройка модулей медиаконвертера \texttt{ansusenko-mc-2}.Мы заменили имеющиеся модули на \texttt{PT-REPEATER-
NM-1FFE} и \texttt{PT-REPEATER-NM-1CFE} (Рис. \ref{05}):

\begin{center}
    \centering
    \includegraphics[width=\textwidth]{../images/05.png}
    \captionof{figure}{Медиаконвертер с доступными модулями расширения.}
    \label{05}
\end{center}

\end{enumerate}

\subsection{Настройка маршрутизатора ansusenko-gw-1}

Маршрутизатор \texttt{ansusenko-gw-1} является шлюзом между локальной сетью организации и сетью провайдера.

\begin{enumerate}
\setcounter{enumi}{4}
\item Конфигурация интерфейсов маршрутизатора \texttt{ansusenko-gw-1}. Конфигурация интерфейса FastEthernet 0/0 с адресом 10.128.0.1/24 (локальная сеть) и интерфейса FastEthernet 0/1 с адресом 192.0.2.1/24 (Рис. \ref{06}):

\begin{center}
    \centering
    \includegraphics[width=\textwidth]{../images/06.png}
    \captionof{figure}{Настройка IP-адресов на маршрутизаторе ansusenko-gw-1.}
    \label{06}
\end{center}

\end{enumerate}

\subsection{Настройка IP-адресов серверов модельного Интернета}

Серверам модельного Интернета назначаются статические IP-адреса из сети 192.0.2.0/24 согласно таблице 11.1.

\begin{enumerate}
\setcounter{enumi}{5}
\item Настройка IP-адреса сервера \texttt{www.yandex.ru} (Рис. \ref{07}):

\begin{center}
    \centering
    \includegraphics[width=\textwidth]{../images/07.png}
    \captionof{figure}{Настройка IP-адреса сервера www.yandex.ru.}
    \label{07}
\end{center}

\texttt{www.yandex.ru}: статический IP-адрес 192.0.2.11, маска 255.255.255.0.}

\item Настройка IP-адреса сервера \texttt{stud.rudn.university} (Рис. \ref{08}):

\begin{center}
    \centering
    \includegraphics[width=\textwidth]{../images/08.png}
    \captionof{figure}{Настройка IP-адреса сервера stud.rudn.university.}
    \label{08}
\end{center}

\item Настройка IP-адреса сервера \texttt{esystem.pfur.ru} (Рис. \ref{09}):

\begin{center}
    \centering
    \centering
    \includegraphics[width=\textwidth]{../images/09.png}
    \captionof{figure}{Настройка IP-адреса сервера esystem.pfur.ru.}
    \label{09}
\end{center}

\item Настройка IP-адреса сервера \texttt{www.rudn.ru} (Рис. \ref{10}):

\begin{center}
    \centering
    \centering
    \includegraphics[width=\textwidth]{../images/10.png}
    \captionof{figure}{Настройка IP-адреса сервера esystem.pfur.ru.}
    \label{10}
\end{center}

\end{enumerate}

\subsection{Настройка DNS-сервера}

На DNS-сервере локальной сети \texttt{ansusenko-dns} добавляются записи для всех серверов модельного Интернета.

\begin{enumerate}
\setcounter{enumi}{8}
\item Конфигурация DNS-сервера \texttt{ansusenko-dns} (Рис. \ref{11}):

\begin{center}
    \centering
    \includegraphics[width=\textwidth]{../images/11.png}
    \captionof{figure}{DNS-сервер с добавленными A-записями.}
    \label{11}
\end{center}

\textit{Таблица DNS-записей:}
\begin{itemize}
    \item \texttt{www.yandex.ru} - 192.0.2.11
    \item \texttt{stud.rudn.university} - 192.0.2.12
    \item \texttt{esystem.pfur.ru} - 192.0.2.13
    \item \texttt{www.rudn.ru} - 192.0.2.14
    \item \texttt{www.ansusenko.rudn.ru} - 10.128.0.2
    \item \texttt{file.ansusenko.rudn.ru} - 10.128.0.3
    \item \texttt{mail.ansusenko.rudn.ru} - 10.128.0.4
    \item \texttt{dns.ansusenko.rudn.ru} - 10.128.0.5
\end{itemize}
\end{enumerate}

\subsection{Схемы L1, L2, L3 и таблицы портов}

Ниже представлены сводные таблицы, отражающие результаты настройки сети.

\subsubsection{Таблица портов подключения оборудования (L1)}

\begin{center}
\begin{tabular}{|l|l|l|l|}
\hline
\textbf{Устройство 1} & \textbf{Порт} & \textbf{Устройство 2} & \textbf{Порт} \\
\hline
msk-ansusenko-gw-1 & fa0/1 & ansusenko-mc-2 & fa0 \\
ansusenko-mc-2 & fa1 & ansusenko-mc-1 & fa1 \\
ansusenko-mc-1 & fa0 & ansusenko-sw-1 & fa0/2 \\
ansusenko-sw-1 & fa0/1 & ansusenko-gw-1 & fa0/0 \\
ansusenko-gw-1 & fa0/1 & ansusenko-mc-02 & fa0 \\
ansusenko-mc-02 & fa1 & ansusenko-mc-01 & fa1 \\
ansusenko-mc-01 & fa0 & internet-sw-1 & fa0/1 \\
internet-sw-1 & fa0/2 & esystem.pfur.ru & fa0 \\
internet-sw-1 & fa0/3 & www.rudn.ru & fa0 \\
internet-sw-1 & fa0/4 & stud.rudn.university & fa0 \\
internet-sw-1 & fa0/5 & www.yandex.ru & fa0 \\
\hline
\end{tabular}
\end{center}

\subsubsection{Таблица IP-адресов (L3)}

\begin{center}
\begin{tabular}{|l|l|l|l|}
\hline
\textbf{Устройство} & \textbf{Интерфейс} & \textbf{IP-адрес} & \textbf{Маска} \\
\hline
ansusenko-gw-1 & fa0/0 & 10.128.0.1 & 255.255.255.0 \\
ansusenko-gw-1 & fa0/1 & 192.0.2.1 & 255.255.255.0 \\
www.yandex.ru & fa0 & 192.0.2.11 & 255.255.255.0 \\
stud.rudn.university & fa0 & 192.0.2.12 & 255.255.255.0 \\
esystem.pfur.ru & fa0 & 192.0.2.13 & 255.255.255.0 \\
www.rudn.ru & fa0 & 192.0.2.14 & 255.255.255.0 \\
ansusenko-dns & — & 10.128.0.5 & 255.255.255.0 \\
\hline
\end{tabular}
\end{center}

\newpage
\section{Выводы}

В ходе выполнения лабораторной работы была построена физическая топология L1, включающая локальную сеть организации, сеть провайдера (медиаконвертеры \texttt{ansusenko-mc-1}, \texttt{ansusenko-mc-2}) и модельную сеть Интернета (коммутатор \texttt{internet-sw-1}, 4 сервера).Настроены IP-адреса на маршрутизаторе \texttt{ansusenko-gw-1}. Подготовлена инфраструктура для последующей настройки механизма NAT.

\end{document}

